\subsection{Two-column homology, reflected cycles, and multiplication}
\label{sec:two-column-homology}

Put $W=W_n$ and $m=\dim W=n-1$. For $0\le q\le p$, let
$\lambda(p,q)=(2^q,1^{p-q})$, the partition with column heights $p,q$.
Its degree is $D=p+q$. Write $M_{p,q}$ for its stable row multiplicity
space in full diagonal harmonics, with $M_{0,0}=K$.
Only the invariant quadratic form survives in this sector; consequently
its multiplicity spaces and moment homology carry an $O(W)$ action
extending their $S_n$ action.

\begin{theorem}[Two-term reduction in every degree]
\label{thm:two-column-complex}
After eliminating the linear moments, the full moment Koszul complex
in row $\lambda(p,q)$ is
\begin{equation}\label{eq:two-column-two-term}
 \Schur_{\lambda(p-1,q-1)}W
                  \xrightarrow{\ L\ }
 \Schur_{\lambda(p,q)}W
\end{equation}
in homological degrees one and zero, with zero source when $q=0$.
The differential inserts the invariant symmetric form.
There is no $H_j$ for $j\ge2$ in any two-column row.
\end{theorem}
\begin{proof}
A moment of order at least three forces first row at least three
in every Littlewood--Richardson constituent. The same is true of
two or more quadratic Koszul generators, since
$\bigwedge^2\Sym^2V=\Schur_{(3,1)}V$ and
$\bigwedge^j\Sym^2V$ is a summand of
$\bigwedge^2\Sym^2V\otimes\bigwedge^{j-2}\Sym^2V$ for $j\ge2$.
Thus only zero or one quadratic generator survives. Cauchy
decomposition gives the target; horizontal Pieri removes one box
from each column to give the source once. The differential is
quadratic multiplication.
\end{proof}

Let $U=K^2$ record two exterior copies of $W$, and put
\[
 \mathcal E=\Lambda(U\otimes W),\qquad
 \omega=\sum_{i=1}^m\theta_i\wedge\phi_i,\qquad
 L=\omega\wedge-,\qquad
 c=\sum_{i=1}^m\partial_{\phi_i}\partial_{\theta_i}.
\]
The displayed orthonormal notation denotes the operators defined by
the form and its inverse, so no choice of an orthonormal basis over
$K$ is required. They satisfy
$[c,L]|_{\mathcal E_e}=(m-e)\operatorname{id}$. For
$P_e=\ker(c:\mathcal E_e\to\mathcal E_{e-2})$, $0\le e\le m$,
the exterior Lefschetz decomposition gives
\begin{align}
 \mathcal E&=\bigoplus_{e=0}^m\ \bigoplus_{j=0}^{m-e}L^jP_e,
 \label{eq:exterior-lefschetz-strings}\\
 cL^jf&=j(m-e-j+1)L^{j-1}f\qquad(f\in P_e).
 \label{eq:exterior-string-normalization}
\end{align}
These follow directly from finite-dimensional $\mathfrak{sl}_2$
theory; see also the exterior Lefschetz treatment in
\cite{KimRhoades}. The operator $L$ commutes with the row
$\mathfrak{sl}_2$ on $U$. On bidegree $(a,b)$, its formula in
$\Lambda^aW\otimes\Lambda^bW$ is
$(-1)^a\sum_i\eps_i\otimes\eps_i$. Exterior multiplication by
$\omega$ fixes these signs uniformly.

Write $\eta$ for the quadratic Koszul generator, of homological degree
one, internal degree two, and $GL(U)$ type $\det U$. The reduced
complex is $\eta\mathcal E\xrightarrow{\,L\,}\mathcal E$, with
$d(\eta f)=Lf$. For its algebra formulas we use the totalized
product, so $\eta$ anticommutes with $U\otimes W$; the complete
sign convention is given in Appendix~\ref{app:odd-derived}.

\begin{theorem}[All two-column homology and its reflection]
\label{thm:two-column-reflection}
For $0\le q\le p\le n$, set $(p^*,q^*)=(n-q,n-p)$.
\begin{enumerate}
\item If $p+q\le n-1$, the sole homology is
      $H_0=M_{p,q}=\Schur_{[\lambda(p,q)]}W$ as $O(W)$ modules.
\item If $p+q=n$, the complex is acyclic.
\item If $p+q\ge n+1$, the sole homology is
      \[
       H_1\cong M_{p^*,q^*}
          \cong M_{p^*,q^*}^{\,*}\otimes(\det W)^2
      \]
      as $O(W)$ modules, hence as $S_n$ modules.
      The reflected row has degree $2n-(p+q)$.
\end{enumerate}
If $p>n$, both terms vanish. These assertions hold in every row
dimension in which the original row occurs; in case~(3),
$p^*\le p-1$, so the reflected row occurs there as well.

More explicitly, if $f\in P_e$ and $k=m-e$, reflection is the cycle map
\begin{equation}\label{eq:exterior-reflected-cycle}
 \rho_e(f)=\frac{\eta L^kf}{k!},
 \qquad
 \rho_e^{-1}(\eta g)=\frac{c^kg}{k!}.
\end{equation}
Before extracting row multiplicity spaces, its source is
$P_e\otimes(\det U)^{k+1}$ and its target has internal degree $2n-e$.
\end{theorem}
\begin{proof}
Each string in \eqref{eq:exterior-lefschetz-strings} contributes its
initial vector to $\mathcal E/L\mathcal E$ and its terminal vector to
$\ker L$. Equation~\eqref{eq:exterior-string-normalization} gives
$L^{k+1}f=0$ and $c^kL^kf=(k!)^2f$, proving
\eqref{eq:exterior-reflected-cycle}.

The row-$\lambda(p,q)$ multiplicity space is the highest-weight kernel
of the transfer
$\Lambda^pW\otimes\Lambda^qW\to
\Lambda^{p+1}W\otimes\Lambda^{q-1}W$.
These kernels fix the Pieri normalization of
\eqref{eq:two-column-two-term}. Its low-degree primitive part is the
orthogonal traceless Schur module: within-column traces vanish by
antisymmetry, and the remaining trace is encoded by $c$.
This is also the admissible case of the orthogonal Littlewood theorem
\cite[Theorem~4.4]{SSW2013}.
If $f$ has $GL(U)$ highest weight $(a,b)$, then $e=a+b$ and
\eqref{eq:exterior-reflected-cycle} has weight
\[
       (a+k+1,b+k+1)=(n-b,n-a).
\]
Exterior Cauchy decomposition transposes this to the stated row
reflection. The strings leave no homology in internal degree $n$.
The dual formula follows from the exterior top-product pairing,
whose orientation after adjoining $\eta$ is
$(\det U)^n\otimes(\det W)^2$.
Finally, if $p>n$, both Schur terms have more than $\dim W$ rows.
\end{proof}

For the bivariate $S_n$-character $R_n(u,v)$ of the exterior $H_0$,
reflection and the Euler character give
\[
 R_n(u,v)-(uv)^nR_n(u^{-1},v^{-1})
    =(1-uv)\lambda_u(W_n)\lambda_v(W_n),
 \qquad
 \lambda_u(W_n)=\sum_j[\Lambda^jW_n]u^j.
\]
Products are tensor products in the representation ring. The dual
operation disappears because $S_n$-modules are self-dual.

\begin{corollary}[The full two-column first syzygies]
\label{cor:two-column-first-syzygies}
In degree $n+1$, for $1\le b\le\lfloor(n+1)/2\rfloor$, the row
$\lambda=(2^b,1^{n+1-2b})$ has
\[
 H_{1,\lambda}
 \cong\Schur_{[(2^{b-1},1^{n+1-2b})]}W_n\big|_{S_n}
 \cong\sgn\otimes\Schur_{[(2^{b-1})]}W_n\big|_{S_n}.
\]
In particular, for $n\ge4$,
\[
 H_{1,(2,2,1^{n-3})}
   \cong S^{(2,1^{n-2})}\oplus S^{(2,2,1^{n-4})}.
\]
\end{corollary}
\begin{proof}
At this boundary the reflected columns are obtained by subtracting
one from each original column. Orthogonal determinant duality
complements the first column in dimension $n-1$, giving the second
formula. For $b=2$ use
$\Schur_{[(2)]}W_n=\Sym^2W_n-K
 \cong S^{(n-1,1)}\oplus S^{(n-2,2)}$ for $n\ge4$.
\end{proof}

Thus at $n=3$ row $(2,2)$ reflects to $(2)$ and has
$H_1\cong S^{(2,1)}$; at $n=4$ row $(2,2,1)$ reflects to $(2,1)$
and has $H_1\cong S^{(2,1,1)}\oplus S^{(2,2)}$.
The exterior $H_0$ character is also obtained in
\cite{KimRhoades,Lentfer}, with explicit double-hook label formulas
in \cite{LentferSign}; the positive homology is the width-two case
of \cite[Theorem~4.4]{SSW2013}. The cycle map above will now identify
its multiplication as well.

\begin{proposition}[Multiplication on reflected homology]
\label{prop:two-column-module}
Put $R=\mathcal E/(\omega)$ and
$I=\eta\operatorname{Ann}_{\mathcal E}(\omega)$.
For $\beta\in U\otimes W$, set
$\iota_\beta=[c,\varepsilon_\beta]$. Then
\begin{equation}\label{eq:two-column-contraction-action}
 \beta\,\rho_e(f)=-\rho_{e-1}(\iota_\beta f),
 \qquad
 I=R\cdot(\eta P_m).
\end{equation}
In particular, the minimal $R$-generator space of $I$ is $\eta P_m$,
all in internal degree $n+1$, and
\[
 \dim(I/R_+I)
   =\binom{2m}{m}-\binom{2m}{m-2}
   =\frac{1}{n+1}\binom{2n}{n}.
\]
The last formula uses the usual zero convention for binomial
coefficients outside their range.
\end{proposition}
\begin{proof}
The identities $[L,\iota_\beta]=\varepsilon_\beta$ and
$[c,\iota_\beta]=0$ give
$\iota_\beta P_e\subseteq P_{e-1}$. With $k=m-e$,
\[
 L^{k+1}\iota_\beta f=(k+1)\beta L^kf
\]
because $L^{k+1}f=0$. Division by $(k+1)!$, followed by moving
$\beta$ across $\eta$, gives the first identity.
The contraction $U\otimes W\otimes P_e\to P_{e-1}$ is a nonzero
symplectic-equivariant map onto an irreducible primitive module.
Contracting $\theta_1\cdots\theta_e$ by $\phi_e$ proves nonvanishing.
It is therefore onto, and iteration from $e=m$ proves generation.
For $m=0$ the assertion is simply $I=K\eta$.
The generator quotient is its first nonzero degree; the dimension
formula is the primitive exterior-power formula.
\end{proof}

The Catalan formula of \cite[Corollary~4.4]{KimRhoades}
counts the minimal syzygy-generating layer after evaluation on two odd rows.
The identities are $Sp_{2m}$-equivariant. For $GL(U)\times O(W)$,
contraction contributes the extra $\det U$ needed between the source
twists of $\rho_e$ and $\rho_{e-1}$; on row multiplicity spaces they
are $O(W)$-equivariant.

\begin{corollary}[All two-column syzygies are generated at the first boundary]
\label{cor:two-column-syzygy-generation}
Write $H_{j,D}(V)$ for the degree-$D$ moment homology as a polynomial
functor of the ordinary row space $V$. For every $D\ge n+1$, the
cokernel of multiplication
\[
 H_{0,D-n-1}(V)\otimes H_{1,n+1}(V)
                         \longrightarrow H_{1,D}(V)
\]
has no constituent $\Schur_\lambda V$ with $\lambda_1\le2$.
This holds in every row dimension in which that constituent occurs.
\end{corollary}
\begin{proof}
Purely odd two-dimensional evaluation is exact in characteristic
zero and sends
$\Schur_\lambda(\Pi U)$ to
$\Pi^{|\lambda|}\Schur_{\lambda'}U$.
It kills precisely the shapes with $\lambda_1>2$ and detects every
remaining Schur multiplicity. It preserves the moment differential
and multiplication. The evaluated map is onto by
\eqref{eq:two-column-contraction-action}, so its original cokernel
has no surviving two-column shape.
\end{proof}

The quotient by shapes with more than two columns carries the induced
algebra structure, since these shapes form a tensor ideal by
Littlewood--Richardson containment.
Appendix~\ref{app:odd-derived} identifies the complete odd differential
algebra and shows that triple Massey products generate its positive
homology.

\begin{proposition}[A single-degree criterion for two-column generation]
\label{prop:two-column-middle-test}
For fixed $n\ge2$, Vandermonde generation in every two-column row and
every degree is equivalent to generation of the two-column rows in
degree $n-1$. It suffices to test the rows
\[
             (2^q,1^{n-1-2q}),
       \qquad 2\le q\le\lfloor(n-1)/2\rfloor.
\]
For $n\ge5$, these tests require only $n-3$ ordinary alphabets.
All two-column rows are generated for $n\le5$.
\end{proposition}
\begin{proof}
Let $G$ be the purely odd two-dimensional evaluation of the ordinary
operator subfunctor $\Pol$. This means evaluation of the entire
operator module, not of its Vandermonde seed alone.
The corresponding full harmonic space is
$\bigoplus_{e=0}^mP_e$.
Derivative closure is expressed by the natural polynomial coaction
\[
 \Pol_D(V)\longrightarrow
            (V\otimes W)\otimes\Pol_{D-1}(V).
\]
Exact odd evaluation therefore makes $G$ closed under every
constant-coefficient exterior contraction.
The map $(U\otimes W)^*\otimes P_e\to P_{e-1}$ is onto, by the
same nonzero symplectic-equivariant contraction used in
Proposition~\ref{prop:two-column-module}.
Thus $G_m=P_m$ propagates down to $G_e=P_e$ for every $e\le m$.
Odd evaluation detects every two-column shape, proving the
equivalence. The cases $q=0,1$ are row hooks and follow from
Proposition~\ref{prop:row-hook-generation}. Each remaining row has
$n-1-q\le n-3$ parts, so Lemma~\ref{lem:stability} and
Corollary~\ref{cor:r3-exact-specialization} make $n-3$ alphabets
detect every remaining multiplicity. The list is empty for $n\le4$;
for $n=5$ its only row is $(2,2)$, covered by Haiman's bivariate
operator theorem as in Theorem~\ref{thm:long-first-row}.
\end{proof}

The remaining task is to generate the displayed degree-$(n-1)$ rows.
The derivative coaction then propagates generation to every lower degree.
